\clearpage
\section*{Public Blog: How Much Comfort Does One Kilowatt-Hour Buy?}
\addcontentsline{toc}{section}{Public blog}

An outdoor fan can help people feel cooler while the air thermometer barely changes. An evaporative cooler can lower air temperature while adding moisture. To compare them fairly, we followed the cooling all the way to the guest: what temperature, humidity, air movement and radiant heat reach that person, and how long do they benefit?

Our model uses Standard Effective Temperature, or SET, to combine those influences into one thermal-response measure. We track relief above a stated reference condition, then divide the total person-time benefit by electricity use. One service unit means one person's SET excess is reduced by one degree for one hour. This is a model-based comfort measure, not a degree of outdoor air cooling or a unit of heat extracted.

Consider our Singapore dining scenario: 120 guests, four hours, and ten carefully placed fans. The calculation gives about \textbf{971 service units for 10 kWh}, or \textbf{97 service units per kWh}. Across all guests, that corresponds to roughly two degrees less SET excess during the event. The fans directly serve about 90\% of guest-time.

The Phoenix example needs a different choice to meet its service target. Ten fans plus two evaporative units supply about \textbf{1,249 service units for 30 kWh}: about \textbf{42 per kWh}. Dry air leaves more room for water to evaporate, but equipment power, delivery losses and where people sit still matter. A higher-powered cooler is worthwhile when its extra relief is needed and reaches enough guests.

One number cannot tell the whole story. Doubling the Singapore dining area and crowd while keeping ten fans leaves the efficiency ratio unchanged, yet halves coverage to 45\%. Check who is served, how much warmth remains, and what the equipment costs. During severe heat, follow local advisories and adjust outdoor plans \cite{nws-heat,nea-heat}.

